% Chương 3

\chapter{PHƯƠNG PHÁP VÀ KIẾN TRÚC HỆ THỐNG}

\label{Chapter3}

%----------------------------------------------------------------------------------------

\section{Kiến trúc hệ thống tổng quát (System Pipeline)}

Hệ thống xử lý khối ảnh y khoa 3D được thiết kế theo mô hình luồng dữ liệu (Data Pipeline), phân chia rõ ràng các tác vụ giữa CPU (Host) và GPU (Device) nhằm tận dụng tối đa băng thông bộ nhớ và năng lực tính toán song song. Cấu trúc hệ thống bao gồm 3 giai đoạn chính:

\begin{enumerate}
    \item \textbf{Giai đoạn 1: Tiền xử lý dữ liệu I/O (CPU):} Chuyển đổi định dạng dữ liệu đầu vào từ dạng văn bản phân mảnh sang dạng nhị phân nguyên khối.
    
    \item \textbf{Giai đoạn 2: Trích xuất bản đồ 2D (GPU):} Nạp toàn bộ khối dữ liệu 3D lên VRAM và thực thi thuật toán MIP để ép xuống mặt phẳng 2D.
    
    \item \textbf{Giai đoạn 3: Xử lý song song đa nhánh (GPU \& CPU):} Từ bản đồ 2D gốc, luồng thực thi được chia làm hai nhánh độc lập:
    \begin{itemize}
        \item \textit{Nhánh 1 (Vessel Analysis):} Xử lý hình thái mạch máu (Làm mờ $\rightarrow$ Phân ngưỡng $\rightarrow$ Đếm điểm ảnh).
        \item \textit{Nhánh 2 (Texture Analysis):} Xử lý kết cấu mô (Khởi tạo GLCM bằng GPU $\rightarrow$ Trích xuất đặc trưng toán học bằng CPU).
    \end{itemize}
\end{enumerate}

% TODO: Vẽ một Sơ đồ khối ba giai đoạn trên

%----------------------------------------------------------------------------------------

\section{Tối ưu hóa I/O Dữ liệu (Data I/O Optimization)}

Một trong những nút thắt cổ chai lớn nhất của hệ thống ban đầu là quá trình nạp dữ liệu. Dữ liệu gốc được lưu dưới dạng hàng trăm tệp văn bản tĩnh (\code{.txt}), chứa các giá trị mức xám dưới dạng chuỗi ký tự ASCII. Việc yêu cầu CPU mở từng tệp, đọc chuỗi, giải mã thành số thực (float) và chuẩn hóa màu sắc mất rất nhiều thời gian.

Để giải quyết vấn đề này, một module Tiền xử lý dữ liệu (Preprocessor) được xây dựng trên C++. Module này thực hiện các bước sau:

\begin{itemize}
    \item Quét qua toàn bộ các tệp \code{.txt}.
    \item Nén và xuất toàn bộ dữ liệu ra một tệp nhị phân duy nhất (\code{volume\_3d.raw}).
\end{itemize}

Nhờ định dạng Binary \code{.raw}, máy tính có thể nạp hàng tỷ byte dữ liệu trực tiếp vào bộ nhớ RAM và đẩy lên VRAM của GPU trong thời gian ngắn.

%----------------------------------------------------------------------------------------

\section{Thiết kế và Triển khai các hàm Kernel CUDA}

Để thực thi song song trên GPU, mạng lưới luồng (Grid) được thiết kế dựa trên kích thước của bề mặt ảnh 2D ($W \times H$). Cụ thể, ảnh được chia thành các Block nhỏ kích thước $16 \times 16$ threads. Số lượng Block được tự động tính toán bù trừ để đảm bảo phủ kín toàn bộ diện tích ảnh mà không bị tràn viền.

\subsection{Kernel Chiếu ảnh cường độ tối đa (MIP\_Projection)}

Thuật toán MIP duyệt qua toàn bộ chiều sâu $D$ của khối ảnh. Trong không gian 3D, mỗi luồng (thread) được gắn tọa độ $(x, y)$ tương ứng với một điểm ảnh trên mặt phẳng đích. GPU phân công mỗi luồng chỉ chạy 1 vòng lặp dọc theo trục $Z$. Luồng sẽ duyệt qua các giá trị tại tọa độ $(x, y, z)$ tương ứng, so sánh và lưu trữ giá trị cực đại tìm được vào biến \code{max\_val}, sau đó ghi thẳng vào mảng \code{d\_mip} 2D.

\subsection{Các Kernel thuộc nhánh Vessel Analysis}

\begin{itemize}
    \item \textbf{Kernel Làm mờ (Mean\_Blur):} Mỗi luồng tại tọa độ $(x, y)$ kiểm tra các điều kiện biên để tránh lỗi truy xuất vùng nhớ (Segmentation fault). Nếu an toàn, luồng sẽ tính tổng cường độ của điểm ảnh đó và 8 điểm lân cận, sau đó chia cho 9 để ra giá trị trung bình và ghi vào ảnh \code{d\_blurred}.
    
    \item \textbf{Kernel Phân ngưỡng (Thresholding):} Các luồng hoạt động hoàn toàn độc lập, thực hiện phép toán rẽ nhánh: Nếu giá trị tại $(x,y)$ vượt qua ngưỡng được cài đặt (ví dụ: 140), gán bằng 255; ngược lại gán bằng 0.
    
    \item \textbf{Kernel Đếm mạch máu (Count\_Vessels):} Để tránh việc hàng ngàn luồng ghi đè lên nhau khi cộng dồn tổng số điểm ảnh trắng, lệnh khóa bộ nhớ \code{atomicAdd()} được sử dụng. Mỗi luồng phát hiện điểm mạch máu sẽ an toàn cộng thêm 1 vào biến con trỏ \code{d\_white\_count} lưu trên VRAM.
\end{itemize}

\subsection{Kernel thuộc nhánh Texture Analysis (GLCM)}

Khác với nhánh trên, tính toán kết cấu yêu cầu đếm tần suất xuất hiện của các cặp điểm ảnh lân cận. Tại hàm \code{Compute\_GLCM}, mỗi luồng xử lý một cặp điểm ảnh (\code{pixel\_1}, \code{pixel\_2}) theo độ dời ngang ($dx = 1, dy = 0$).

Ma trận GLCM $256 \times 256$ được trải phẳng thành mảng 1 chiều kích thước 65.536 phần tử. Chỉ số để cập nhật ma trận được tính toán bằng công thức:
\begin{equation}
    \text{index} = \text{pixel\_1} \times 256 + \text{pixel\_2}
    \label{eq:glcm_index}
\end{equation}

Tại chỉ số này, luồng tiếp tục gọi \code{atomicAdd()} để tăng tần suất cặp mức xám lên 1 đơn vị.

%----------------------------------------------------------------------------------------

\section{Quản lý bộ nhớ và Đồng bộ hóa}

\begin{itemize}
    \item \textbf{Cấp phát:} Sử dụng \code{cudaMalloc} để cấp phát vùng nhớ độc lập trên VRAM cho từng giai đoạn (Khối gốc, Ảnh MIP, Ảnh Blur, Ảnh nhị phân và Ma trận GLCM).
    
    \item \textbf{Làm sạch bộ nhớ:} Trước khi khởi chạy các Kernel đếm tổng hoặc GLCM, bộ nhớ VRAM luôn chứa dữ liệu rác. Hàm \code{cudaMemset} được gọi để ép các vùng nhớ này về 0, đảm bảo tính chính xác tuyệt đối cho \code{atomicAdd}.
    
    \item \textbf{Đồng bộ hóa:} Giữa mỗi bước thực thi Kernel, lệnh \code{cudaDeviceSynchronize()} đóng vai trò là rào chắn (Barrier), bắt buộc CPU phải tạm dừng để chờ toàn bộ hàng chục ngàn luồng của GPU hoàn thành xong công việc hiện tại, mới được phép phát lệnh chạy Kernel tiếp theo, tránh tình trạng lấy kết quả sai lệch.
\end{itemize}

%----------------------------------------------------------------------------------------
